\section{Further research and shortcuts that fail}
The improvements leave several concrete research directions. None is asserted as a theorem or an optimality result here.
\begin{enumerate}
\item Seek a lower-degree carrier with the same no-rational-root and every-modulus-root properties while preserving the quantitative powered-projection bound. Reducing the formal carrier degree alone is insufficient if it increases the strong-scale or norm degree.
\item Investigate whether additional common prime factors in the auxiliary modulus can be absorbed into a supported-modulus interpolation family. The strong pair must remain fixed after a residue is prescribed, and the full squared-modulus period must be proved.
\item Search the Pareto frontier of selective even-index choices and residue carriers. Doubling all indices is not a harmless uniform simplification: the retained ordinary $t$ seed is what makes the one-step $t+1$ addition valid and helps bootstrap the radix.
\item Optimize literal monic projection and signed factorization jointly, while requiring full coefficient reconstruction and explicit inverse substitutions. A smaller residual count alone does not guarantee fewer variables, terms, or a lower height.
\item Formalize the new unbounded interfaces in a proof assistant, keeping the already named native and C222 calibration obligations visible. This would be a separate deliverable from the exact finite algebra checks supplied here.
\end{enumerate}

Several tempting shortcuts are invalid. A square product does not imply that both factors are squares without coprimality and positivity. A powered congruence $H\mid K^5$ does not imply $H\mid K$; the strict size argument is essential. A rational ordinate of an adjacent Pell equation cannot simply be declared integral; descent is required. The original period $4m_0$ does not generally suffice modulo $f^2$; the corrected factor $f$ must be retained. The square-output factorial value cannot be assumed before proving radix carry bounds and the ceiling fixed point. Replacing a remainder by $1+qu$ is not a bijective change of arbitrary integer coordinates. Residue one cannot orient a signed power when $q=1$ or $2$. Finally, changing a seed index does not preserve its old private Pell witnesses; existence must be proved by reconstructing them.

The supplied family already includes a quartic in $37$ unrestricted integers with enormous coefficients, and ten-integer equations using a cited Sun reduction. Accordingly, neither $442$ variables for the direct quartic nor eleven variables for the norm is advertised as a global variable-count record. The purpose is a transparent, fully specified and reproducible improvement in the stated multidimensional objectives.
